\documentclass[11pt]{article}
\usepackage[margin=1in]{geometry}
\usepackage{amsmath,amssymb,amsthm}
\usepackage{hyperref}
\newtheorem{theorem}{Theorem}
\newtheorem{lemma}[theorem]{Lemma}
\newtheorem{proposition}[theorem]{Proposition}
\theoremstyle{definition}
\newtheorem{definition}[theorem]{Definition}
\theoremstyle{remark}
\newtheorem{remark}[theorem]{Remark}

\title{Disproof of Conjecture 00000001028:\\
there are not 21 exceptional values below 316 for Kirkman quadruple systems}
\author{jilint777}
\date{2026-10-04}

\begin{document}
\sloppy
\maketitle

\begin{abstract}
Conjecture 00000001028 asserts that Kirkman quadruple systems $\mathrm{KQS}(v)$
(resolvable $S(2,4,v)$ designs) exist for all $v\equiv4\pmod{12}$ with $v\ge v_0=316$, and
that below $v_0$ there are exactly $21$ exceptional values. Only $26$ values
$v\equiv 4\pmod{12}$ lie below $316$. We construct explicit Kirkman quadruple systems for
the seven values $v=4,16,28,40,52,64,76$. Hence at most $19$ values below $316$ are
exceptional (at most $20$ even if $v_0=316$ itself is counted), and the second clause of the
conjecture is false. The seven designs, the bound and the negation of the conjecture are
formalized and machine-checked in Lean~4 (core library only). (In fact, by a 1972 theorem of
Hanani, Ray-Chaudhuri and Wilson, there are no exceptional values at all.)
\end{abstract}

\section{The conjecture and its reading}
\begin{quote}
\emph{Definition:} A resolvable $S(2,4,v)$ (Kirkman quadruple system) is a design whose
blocks partition into parallel classes. \emph{Conjecture:} There is an explicit threshold
$v_0=316$ such that $\mathrm{KQS}(2,4,v)$ exist for all $v\equiv 4\bmod 12$ with $v\ge v_0$;
below $v_0$ there are exactly $21$ exceptional values.
\end{quote}
The Chinese version says the same. The conjecture is a conjunction of
\begin{itemize}
\item[(C1)] $\mathrm{KQS}(v)$ exists for every $v\equiv4\pmod{12}$ with $v\ge316$, and
\item[(C2)] below $v_0=316$ there are exactly $21$ exceptional values.
\end{itemize}
We refute (C2); this refutes the conjecture regardless of (C1).

\paragraph{Exceptional values.} In the context of (C1), an \emph{exceptional value} is an
admissible order $v\equiv4\pmod{12}$, $v<v_0$, for which no $\mathrm{KQS}(v)$ exists.
(C2) says that the set of these values has exactly $21$ elements. We show that it has at
most $19$ elements. We also cover the following variants.
\begin{itemize}
\item \emph{Non-strict ``below''} ($v\le316$): at most $20$ exceptional values, so still not
$21$. (Under (C1), $v=316$ is not exceptional anyway.)
\item \emph{Excluding trivial orders} (for example, ignoring $v=4$): the set only gets
smaller, so the bound $\le 19$ still holds.
\item \emph{Any residue:} if ``exceptional'' meant any $v<316$ without a KQS, regardless of
the residue of $v$, then the count would be at least $288$ (Remark~\ref{rem:nec}), again not
$21$. This variant is discussed only in this report; the Lean development treats the
main reading and the non-strict variant.
\item \emph{``Exceptional'' as ``existence undecided'':} this is not a property of $v$ alone.
In any case, the seven values below are settled by explicit, verified designs, so at most
$19$ values could be undecided. By \cite{hrw}, none is.
\end{itemize}

\section{Definitions}
\begin{definition}
Let $v\ge1$ and $V=\{0,1,\dots,v-1\}$. A \emph{Kirkman quadruple system}
$\mathrm{KQS}(v)$ is a family $\mathcal R=(P_1,\dots,P_r)$ of \emph{parallel classes}. Each
$P_i$ is a family of \emph{blocks}, and the following hold:
\begin{enumerate}
\item every block is a $4$-element subset of $V$;
\item every parallel class partitions $V$: each point of $V$ lies in exactly one block of
$P_i$;
\item every pair of distinct points of $V$ lies in exactly one block of
$P_1\cup\dots\cup P_r$.
\end{enumerate}
Conditions 1 and 3 say that the blocks form an $S(2,4,v)$, i.e.\ a $2$-$(v,4,1)$ design;
condition 2 says that $\mathcal R$ is a resolution of it.
\end{definition}

\begin{remark}\label{rem:nec}
If a $\mathrm{KQS}(v)$ exists and $v\ge2$, then $v\equiv4\pmod{12}$. Indeed, the blocks through a
fixed point $x$ partition $V\setminus\{x\}$ into triples, so $3\mid v-1$. A parallel class
partitions $V$ into $4$-sets, so $4\mid v$. Hence there are exactly $26$ admissible values
below $316$, namely $4,16,28,\dots,304$, and $27$ up to $316$. Moreover, every $v<316$ with
$2\le v$ and $v\not\equiv4\pmod{12}$ has no KQS. There are $314-26=288$ such values, which is
why the ``any residue'' reading gives at least $288$ exceptions.
\end{remark}

\section{Seven Kirkman quadruple systems}
\label{sec:constr}

\subsection{\texorpdfstring{$v=4$}{v=4}} One parallel class consisting of the single block
$\{0,1,2,3\}$.

\subsection{\texorpdfstring{$v=16$ and $v=64$}{v=16 and v=64}: affine spaces over GF(4)}
Let $\mathbb F_4=\{0,1,\omega,\omega^2\}$, encoded as $\{0,1,2,3\}$; addition is bitwise XOR and
$\omega\cdot\omega=\omega^2$. The points are the vectors of $\mathbb F_4^n$ ($n=2,3$),
and $(a_0,\dots,a_{n-1})$ is encoded as $\sum a_i4^i$. The blocks are the affine lines
$\{p+td:t\in\mathbb F_4\}$ with $d\neq0$. Two distinct points $p,q$ lie on exactly one line,
namely the one with direction $q-p$, since the direction of a line through $p$ and $q$ is a
nonzero multiple of $q-p$. For each direction $d$ (up to scalars there are $(4^n-1)/3$ of them),
the lines with direction $d$ are the cosets of $\mathbb F_4d$ and partition the space. So the
lines of $\mathrm{AG}(n,4)$ form a $\mathrm{KQS}(4^n)$ with $(4^n-1)/3$ parallel classes:
$5$ classes of $4$ blocks for $v=16$ and $21$ classes of $16$ blocks for $v=64$.

\subsection{\texorpdfstring{$v=28$}{v=28}: the Hermitian unital}
Let $\mathbb F_9=\mathbb F_3[i]$ with $i^2=-1$. The Hermitian curve $x^4+y^4+z^4=0$ in
$\mathrm{PG}(2,9)$ has $28$ points. Every line of $\mathrm{PG}(2,9)$ meets it in $1$ or $4$
points, so the $63$ secant lines form a $2$-$(28,4,1)$ design, the Hermitian unital $U(3)$.
An exact-cover search found a resolution of it into $9$ parallel classes of $7$ blocks.
The labelling is as follows. We write $a+bi$ as the pair $(a,b)\in\{0,1,2\}^2$, normalize
projective points so that the first nonzero coordinate is $1$, and number the $28$ unital
points $0,\dots,27$ in lexicographic order of their coordinate triples. The parallel classes are:

\begin{center}\small
\begin{tabular}{@{}ll@{}}
$P_{1}$: & $\{0,13,16,26\}$ $\{1,5,11,12\}$ $\{2,4,17,25\}$ $\{3,9,21,22\}$ $\{6,8,19,20\}$ $\{7,15,23,27\}$ $\{10,14,18,24\}$\\
$P_{2}$: & $\{0,7,10,12\}$ $\{1,13,20,22\}$ $\{2,9,16,27\}$ $\{3,6,18,23\}$ $\{4,8,24,26\}$ $\{5,14,17,21\}$ $\{11,15,19,25\}$\\
$P_{3}$: & $\{0,5,19,22\}$ $\{1,8,17,27\}$ $\{2,6,10,13\}$ $\{3,15,20,24\}$ $\{4,12,18,21\}$ $\{7,9,25,26\}$ $\{11,14,16,23\}$\\
$P_{4}$: & $\{0,1,2,3\}$ $\{4,14,22,27\}$ $\{5,10,23,26\}$ $\{6,15,16,21\}$ $\{7,11,17,20\}$ $\{8,13,18,25\}$ $\{9,12,19,24\}$\\
$P_{5}$: & $\{0,8,21,23\}$ $\{1,7,16,24\}$ $\{2,14,19,26\}$ $\{3,4,11,13\}$ $\{5,9,18,20\}$ $\{6,12,25,27\}$ $\{10,15,17,22\}$\\
$P_{6}$: & $\{0,6,17,24\}$ $\{1,4,19,23\}$ $\{2,7,18,22\}$ $\{3,5,16,25\}$ $\{8,9,10,11\}$ $\{12,13,14,15\}$ $\{20,21,26,27\}$\\
$P_{7}$: & $\{0,11,18,27\}$ $\{1,6,9,14\}$ $\{2,5,8,15\}$ $\{3,12,17,26\}$ $\{4,10,16,20\}$ $\{7,13,19,21\}$ $\{22,23,24,25\}$\\
$P_{8}$: & $\{0,4,9,15\}$ $\{1,10,21,25\}$ $\{2,12,20,23\}$ $\{3,7,8,14\}$ $\{5,13,24,27\}$ $\{6,11,22,26\}$ $\{16,17,18,19\}$\\
$P_{9}$: & $\{0,14,20,25\}$ $\{1,15,18,26\}$ $\{2,11,21,24\}$ $\{3,10,19,27\}$ $\{4,5,6,7\}$ $\{8,12,16,22\}$ $\{9,13,17,23\}$\\
\end{tabular}
\end{center}

\subsection{\texorpdfstring{$v=40$}{v=40}: a packing of PG(3,3)}
The $40$ points and $130$ lines of $\mathrm{PG}(3,3)$ form a $2$-$(40,4,1)$ design (two
points span one line). A \emph{spread} is a set of $10$ pairwise disjoint lines; a
\emph{packing} is a partition of the $130$ lines into $13$ spreads, i.e.\ a resolution.
Packings of $\mathrm{PG}(3,q)$ exist for all $q$ (Denniston \cite{den}). An exact-cover search
gave the following one. Points are the vectors of $\{0,1,2\}^4\setminus\{0\}$ with first
nonzero coordinate $1$, numbered $0,\dots,39$ in lexicographic order.

\begin{center}\small
\begin{tabular}{@{}ll@{}}
$P_{1}$: & $\{0,34,35,36\}$ $\{1,6,9,12\}$ $\{2,15,16,20\}$ $\{3,22,27,29\}$ $\{4,21,30,39\}$\\
& $\{5,13,23,33\}$ $\{7,14,26,38\}$ $\{8,18,28,32\}$ $\{10,19,25,31\}$ $\{11,17,24,37\}$\\[2pt]
$P_{2}$: & $\{0,25,26,27\}$ $\{1,5,8,11\}$ $\{2,32,36,37\}$ $\{3,15,17,19\}$ $\{4,13,22,31\}$\\
& $\{6,20,28,39\}$ $\{7,18,30,33\}$ $\{9,21,23,34\}$ $\{10,14,29,35\}$ $\{12,16,24,38\}$\\[2pt]
$P_{3}$: & $\{0,31,32,33\}$ $\{1,22,25,28\}$ $\{2,13,17,21\}$ $\{3,5,7,12\}$ $\{4,20,29,38\}$\\
& $\{6,18,26,34\}$ $\{8,14,27,37\}$ $\{9,19,24,35\}$ $\{10,15,30,36\}$ $\{11,16,23,39\}$\\[2pt]
$P_{4}$: & $\{0,19,20,21\}$ $\{1,32,35,38\}$ $\{2,22,26,30\}$ $\{3,4,9,11\}$ $\{5,14,24,31\}$\\
& $\{6,17,25,36\}$ $\{7,15,27,39\}$ $\{8,16,29,33\}$ $\{10,13,28,34\}$ $\{12,18,23,37\}$\\[2pt]
$P_{5}$: & $\{0,7,8,9\}$ $\{1,31,34,37\}$ $\{2,23,27,28\}$ $\{3,14,16,21\}$ $\{4,17,26,35\}$\\
& $\{5,15,22,32\}$ $\{6,19,30,38\}$ $\{10,18,24,39\}$ $\{11,13,29,36\}$ $\{12,20,25,33\}$\\[2pt]
$P_{6}$: & $\{0,37,38,39\}$ $\{1,13,16,19\}$ $\{2,6,7,11\}$ $\{3,23,25,30\}$ $\{4,15,24,33\}$\\
& $\{5,17,27,34\}$ $\{8,21,22,35\}$ $\{9,18,29,31\}$ $\{10,20,26,32\}$ $\{12,14,28,36\}$\\[2pt]
$P_{7}$: & $\{0,1,2,3\}$ $\{4,18,27,36\}$ $\{5,19,29,39\}$ $\{6,13,24,32\}$ $\{7,20,23,35\}$\\
& $\{8,15,25,38\}$ $\{9,17,28,33\}$ $\{10,16,22,37\}$ $\{11,14,30,34\}$ $\{12,21,26,31\}$\\[2pt]
$P_{8}$: & $\{0,10,11,12\}$ $\{1,33,36,39\}$ $\{2,24,25,29\}$ $\{3,13,18,20\}$ $\{4,14,23,32\}$\\
& $\{5,21,28,38\}$ $\{6,16,27,35\}$ $\{7,19,22,34\}$ $\{8,17,30,31\}$ $\{9,15,26,37\}$\\[2pt]
$P_{9}$: & $\{0,13,14,15\}$ $\{1,24,27,30\}$ $\{2,4,8,12\}$ $\{3,32,34,39\}$ $\{5,18,25,35\}$\\
& $\{6,21,29,37\}$ $\{7,16,28,31\}$ $\{9,20,22,36\}$ $\{10,17,23,38\}$ $\{11,19,26,33\}$\\[2pt]
$P_{10}$: & $\{0,22,23,24\}$ $\{1,15,18,21\}$ $\{2,33,34,38\}$ $\{3,6,8,10\}$ $\{4,19,28,37\}$\\
& $\{5,16,26,36\}$ $\{7,17,29,32\}$ $\{9,14,25,39\}$ $\{11,20,27,31\}$ $\{12,13,30,35\}$\\[2pt]
$P_{11}$: & $\{0,4,5,6\}$ $\{1,23,26,29\}$ $\{2,14,18,19\}$ $\{3,31,36,38\}$ $\{7,13,25,37\}$\\
& $\{8,20,24,34\}$ $\{9,16,30,32\}$ $\{10,21,27,33\}$ $\{11,15,28,35\}$ $\{12,17,22,39\}$\\[2pt]
$P_{12}$: & $\{0,28,29,30\}$ $\{1,14,17,20\}$ $\{2,5,9,10\}$ $\{3,33,35,37\}$ $\{4,16,25,34\}$\\
& $\{6,15,23,31\}$ $\{7,21,24,36\}$ $\{8,13,26,39\}$ $\{11,18,22,38\}$ $\{12,19,27,32\}$\\[2pt]
$P_{13}$: & $\{0,16,17,18\}$ $\{1,4,7,10\}$ $\{2,31,35,39\}$ $\{3,24,26,28\}$ $\{5,20,30,37\}$\\
& $\{6,14,22,33\}$ $\{8,19,23,36\}$ $\{9,13,27,38\}$ $\{11,21,25,32\}$ $\{12,15,29,34\}$\\[2pt]
\end{tabular}
\end{center}

\subsection{\texorpdfstring{$v=52$ and $v=76$}{v=52 and v=76}: 1-rotational designs}
\begin{lemma}\label{lem:rot}
Let $m\ge2$, $N=3m$, and let $B_1,\dots,B_s\subset\mathbb Z_N$ be $4$-sets ($4s=m-1$) such that
\begin{enumerate}
\item[(a)] the $4s$ elements of $B_1\cup\dots\cup B_s$, reduced mod $m$, are exactly
$1,2,\dots,m-1$, each once;
\item[(b)] the $12s$ differences $x-y$ ($x\ne y$ in the same $B_i$) are exactly the elements of
$\mathbb Z_N\setminus\{0,m,2m\}$, each once.
\end{enumerate}
On $V=\mathbb Z_N\cup\{\infty\}$, for $j=0,\dots,m-1$ let
$P_j=\bigl\{\{j,j+m,j+2m,\infty\}\bigr\}\cup\{B_i+j+mk: 1\le i\le s,\ 0\le k\le2\}$.
Then $(P_0,\dots,P_{m-1})$ is a $\mathrm{KQS}(3m+1)$.
\end{lemma}
\begin{proof}
\emph{Classes.} By (a), the sets $B_i+mk$ ($k=0,1,2$) cover each residue class $r\not\equiv0\pmod m$
of $\mathbb Z_N$ exactly once (the element $x\equiv r$ yields $x,x+m,x+2m$). After shifting by
$j$, they cover every residue class except that of $j$, and $\{j,j+m,j+2m,\infty\}$ covers the rest.
\emph{Pairs.} The pair $\{\infty,y\}$ lies only in $\{j,j+m,j+2m,\infty\}$ with $j\equiv y\pmod m$.
Let $y\ne z$ be finite and $d=z-y$. If $d\in\{m,2m\}$, then $y\equiv z\pmod m$; by (a), no
translate of a $B_i$ contains two elements of the same residue class mod $m$, so the only block
containing $y,z$ is the $\infty$-block of $j\equiv y$. Otherwise, by (b), $d=x-x'$ for a unique
ordered pair $(x',x)$ in a unique $B_i$. A translate $B_i+t$ contains $y,z$ exactly when
$t=y-x'$, and no other $B_{i'}+t'$ contains them, by uniqueness in (b). Finally, the translates
$B_i+j+mk$ ($0\le j<m$, $0\le k<3$) are the $N$ translates $B_i+t$, $t\in\mathbb Z_N$, each
once. Thus every pair lies in exactly one block.
\end{proof}
A backtracking search gave base blocks satisfying (a) and (b).
\begin{itemize}
\item $m=17$, $v=52$ ($\infty=51$): $\{1,2,20,24\},\{4,25,12,49\},\{5,10,45,30\},\{6,9,48,50\}$.
\item $m=25$, $v=76$ ($\infty=75$): $\{1,34,47,48\},\{2,33,67,69\},\{3,56,7,71\},\{4,61,13,16\},$
$\{5,60,37,43\},\{14,65,70,49\}$.
\end{itemize}
For instance, for $m=17$ the residues mod $17$ of the four base blocks are
\[\{1,2,3,7\},\ \{4,8,12,15\},\ \{5,10,11,13\},\ \{6,9,14,16\}.\]
Conditions (a) and (b) are checked by
\texttt{verify.py}. Independently, the Lean development checks the resulting designs directly
($17$ classes and $221$ blocks for $v=52$; $25$ classes and $475$ blocks for $v=76$).

\section{The disproof}
\begin{theorem}\label{thm:main}
There are at most $19$ values $v\equiv4\pmod{12}$, $v<316$, without a $\mathrm{KQS}(v)$, and
at most $20$ with $v\le316$. Hence clause (C2) and Conjecture 00000001028 are false.
\end{theorem}
\begin{proof}
The admissible values below $316$ are $4+12k$, $0\le k\le25$, i.e.\ $26$ values. By
Section~\ref{sec:constr}, $\mathrm{KQS}(v)$ exists for $v\in\{4,16,28,40,52,64,76\}$. So every
exceptional value lies in the remaining set
$\{88,100,\dots,304\}$ of $19$ values (together with $316$ in the non-strict reading: $20$
values). A set of exactly $21$ exceptional values would have to be contained in a set with at
most $20$ elements, which is impossible.
\end{proof}

\begin{remark}
The truth is much stronger. Hanani, Ray-Chaudhuri and Wilson \cite{hrw} proved that a
$\mathrm{KQS}(v)$ exists for \emph{every} $v\equiv4\pmod{12}$, so there are no exceptional
values at all. Our disproof does not depend on this theorem. It uses only the seven explicit,
machine-verified designs above.
\end{remark}

\section{Lean formalization}
The directory \texttt{lean4/} is a self-contained Lean~4.19.0 project (core library only, no
Mathlib). The main file is \texttt{Main.lean}.
\begin{itemize}
\item \texttt{Resolution} is a list of parallel classes; a class is a list of blocks; a block
is a list of natural numbers. \texttt{IsKQS v R} states conditions 1--3 of the definition
literally: every block has length $4$, has no duplicates and lies in $\{0,\dots,v-1\}$; for
each class $C$ and each $x<v$, the number of blocks of $C$ containing $x$ is $1$; and for all
$x\ne y$ below $v$, the number of blocks of $R$ containing both is $1$.
\texttt{HasKQS v} is $\exists R,\ \texttt{IsKQS v R}$.
\item \texttt{Exceptional v} is $v\bmod 12=4\wedge v<316\wedge\neg\texttt{HasKQS v}$.
\texttt{ExactlyExceptional k} says that the exceptional values are exactly the entries of a
duplicate-free list of length $k$. \texttt{Clause1}, \texttt{Clause2} (that is,
\texttt{ExactlyExceptional 21}) and \texttt{Conjecture} $=$ \texttt{Clause1} $\wedge$
\texttt{Clause2} formalize the statement. \texttt{ExceptionalLe} and
\texttt{ExactlyExceptionalLe} formalize the non-strict variant.
\item \texttt{check v R} is a fast Boolean checker that uses bitmasks
($\sum_{a\in B}2^a$). For each class, the block masks must be pairwise disjoint and cover
$\{0,\dots,v-1\}$. For each point $x$, the masks of the blocks through $x$, with bit $x$
cleared, must be pairwise disjoint and cover all $y\ne x$. The theorem
\texttt{check\_sound : check v R = true -> IsKQS v R} is proved in general, via the counting
lemma \texttt{disjUnion\_sound}: if a disjointness-checked union succeeds, then the number of
masks having bit $y$, plus bit $y$ of the start value, equals bit $y$ of the result.
\item \texttt{R4}, \dots, \texttt{R76} are the seven resolutions, written out as explicit lists.
\texttt{kqs\_4}, \dots, \texttt{kqs\_76} prove \texttt{IsKQS v Rv} by
\texttt{check\_sound} and kernel evaluation (\texttt{decide +kernel}).
\item \texttt{exceptional\_le\_19}: any duplicate-free list of exceptional values has length at
most $19$. \texttt{not\_exactlyExceptional}: no $k\ge20$ is the exact count.
\texttt{clause2\_false : }$\neg$\texttt{Clause2}, and
\texttt{conjecture\_00000001028\_false : }$\neg$\texttt{Conjecture}.
\texttt{exceptionalLe\_le\_20} and \texttt{clause2\_le\_false} handle the non-strict reading.
\texttt{admissible\_below\_316} lists the $26$ admissible values below $316$.
\item Non-vacuity: \texttt{not\_hasKQS\_two\_three} and \texttt{not\_hasKQS\_five} show that
\texttt{HasKQS} is not trivially true. There is no KQS on $2$, $3$ or $5$ points. For $v=5$ a
block exists, but it misses a point $y$, and the block through $y$ in the same parallel class
meets the first block, which contradicts the partition property. Also, \texttt{exceptional\_iff} shows that for
admissible $v<316$, \texttt{Exceptional v} is equivalent to $\neg$\texttt{HasKQS v}.
\end{itemize}
\texttt{lake build} succeeds with no errors or warnings (about one minute of kernel
evaluation). There is no \texttt{sorry}, no \texttt{native\_decide}, and no added axiom;
\texttt{\#print axioms} shows at most \texttt{propext}, \texttt{Classical.choice} and
\texttt{Quot.sound}.

\paragraph{What is not in Lean.} The geometric descriptions of the designs (affine spaces,
unital, packing, Lemma~\ref{lem:rot}) explain where the designs come from. They are not needed
for the proof, because each design is verified directly as an explicit object. The necessary
condition of Remark~\ref{rem:nec} (used only for the ``any residue'' reading) and the
Hanani--Ray-Chaudhuri--Wilson theorem are not formalized.

\section{Independent check}
\texttt{verify.py} (Python~3 standard library) reads the seven resolutions from
\texttt{Main.lean}. It checks them by explicit pair counting, which is a different method from the Lean
bitmask checker. It also re-derives each design from its description: lines of
$\mathrm{AG}(2,4)$ and $\mathrm{AG}(3,4)$, secant lines of the Hermitian curve over
$\mathbb F_9$, lines of $\mathrm{PG}(3,3)$, and the developments of Lemma~\ref{lem:rot},
including conditions (a) and (b). Finally, it recounts the admissible values and the bound.

\section{Reproduction}
\begin{verbatim}
cd lean4 && lake build     # Lean 4.19.0, core only
python3 verify.py          # Python 3 standard library
pdflatex report.tex
\end{verbatim}

\begin{thebibliography}{9}
\bibitem{hrw} H.~Hanani, D.~K.~Ray-Chaudhuri and R.~M.~Wilson, On resolvable designs,
\emph{Discrete Math.} 3 (1972), 343--357.
\bibitem{den} R.~H.~F.~Denniston, Some packings of projective spaces,
\emph{Atti Accad. Naz. Lincei Rend. Cl. Sci. Fis. Mat. Natur.} 52 (1972), 36--40.
\bibitem{crc} C.~J.~Colbourn and J.~H.~Dinitz (eds.), \emph{Handbook of Combinatorial
Designs}, 2nd ed., Chapman \& Hall/CRC, 2007 (resolvable designs, unitals, 1-rotational
constructions).
\end{thebibliography}
\end{document}
